\section{Tujuan Buku Ajar}

Buku ajar ini dirancang sebagai panduan komprehensif untuk menguasai Pemrograman Game menggunakan Unity dan bahasa C\# \cite{unityManual}. Fokus utama buku ini adalah implementasi praktis melalui proyek ``Roll a Ball'' yang saling terhubung dari bab ke bab. Tujuan spesifik buku ini adalah:

\begin{enumerate}
  \item Memberikan pemahaman mendalam tentang konsep pengembangan game dan arsitektur game engine
  \item Mengembangkan kemampuan mengoperasikan Unity Editor dan mengatur Game Objects
  \item Membangun keterampilan scripting C\# untuk perilaku game (input, fisika, collision)
  \item Memfasilitasi pembangunan game 3D lengkap (Roll a Ball) dari awal hingga build
  \item Menumbuhkan kemampuan membuat UI, mengelola game state, dan melakukan build/export
  \item Memfasilitasi pencapaian CPL dan CPMK yang telah ditetapkan
\end{enumerate}

Setelah mempelajari buku ini secara menyeluruh, mahasiswa diharapkan mampu:
\begin{itemize}
  \item Menjelaskan peran game engine dalam pengembangan game
  \item Menggunakan Unity Editor (Hierarchy, Inspector, Scene, Game view)
  \item Menulis script C\# dengan MonoBehaviour (Update, FixedUpdate)
  \item Mengimplementasikan input player, kontrol kamera, dan collision detection
  \item Membuat game Roll a Ball lengkap dengan collectibles, skor, dan UI
  \item Melakukan build game untuk platform target (Windows/standalone)
\end{itemize}
